\documentclass[11pt]{article}
\usepackage[T1]{fontenc}
\usepackage{lmodern}
\usepackage[margin=1in]{geometry}
\usepackage{amsmath,amssymb,amsthm,mathtools,mathrsfs}
\usepackage{microtype,booktabs,array,longtable}
\usepackage[hidelinks]{hyperref}
\hypersetup{pdftitle={RA-04: Clustered-gap bounds for randomized block Krylov approximation},pdfauthor={Sidney Holden},pdfsubject={Full proof; independent informal AI-agent audit passed}}
\usepackage{enumitem}
\usepackage{fancyvrb}
\setlength{\parindent}{0pt}
\setlength{\parskip}{5pt plus 1pt minus 1pt}
\setlength{\emergencystretch}{2em}
\numberwithin{equation}{section}
\newtheorem{theorem}{Theorem}[section]
\newtheorem{lemma}[theorem]{Lemma}
\newtheorem{proposition}[theorem]{Proposition}
\newtheorem{corollary}[theorem]{Corollary}
\theoremstyle{definition}
\newtheorem{definition}[theorem]{Definition}
\newtheorem{remark}[theorem]{Remark}
\newcommand{\R}{\mathbb R}
\newcommand{\C}{\mathbb C}
\newcommand{\K}{\mathcal K}
\newcommand{\E}{\mathbb E}
\newcommand{\Prob}{\mathbb P}
\newcommand{\diag}{\operatorname{diag}}
\newcommand{\rank}{\operatorname{rank}}
\newcommand{\dist}{\operatorname{dist}}
\newcommand{\range}{\operatorname{range}}
\newcommand{\sm}{\sigma_{\min}}
\newcommand{\nn}[1]{\lVert #1\rVert}
\newcommand{\ip}[2]{\langle #1,#2\rangle}
\newcommand{\ahat}{\widehat A}
\newcommand{\Bconst}{128e}
\title{RA-04: Clustered-gap bounds for randomized block Krylov approximation}
\author{Sidney Holden\\{\small Center for Computational Biology, Flatiron Institute}\\{\small Simons Foundation}}
\date{}
\begin{document}
\maketitle
\vspace{-1em}
\begin{abstract}
A complete proof of the displayed RA-04 assertion is claimed. For a positive head spectrum of size $m=bt$ whose relative $b$-step gap is at least $\Delta$, the tangential interpolation evaluation matrix satisfies
\[
\Prob\!\left\{\sup_{0\le x\le\lambda_m}\nn{E_{\lambda,H}(x)}_2
\le 2^{24}m^8\eta^{-4}\left(\frac{128e}{\Delta}\right)^t\right\}\ge1-\eta.
\]
The proof combines a logarithmic root covering, an annular resolvent identity, and a variational descent estimate. A published real-algebraic tube theorem controls proximity to the singular head-Krylov locus. Its defining polynomial has degree $m$, independently of the spectral scale. Translation preserves that locus, so one probability event controls every tail evaluation. The standard good-start convergence theorem then gives the exact iteration order requested in RA-04, without an additional $t\log m$ term.
\end{abstract}

\textbf{Review status.} The full exact RA-04 target passed a separate Codex AI-agent informal mathematical audit on 12 September 2026 (America/New\_York). The repository submission proposes \emph{Solved} under its independent-audit policy. This is not external human peer review or formal verification. No Lean verification was performed. The independent report is included in \texttt{verification/independent-review.md}. Finite tests are diagnostics, not proof certificates.

\section{The assertion and the new estimate}\label{sec:statement}
Let $A\in\R^{n\times d}$ and $M=AA^T$. Order the eigenvalues of $M$ as
$\lambda_1\ge\cdots\ge\lambda_n\ge0$, extending this sequence by zeros for indices greater than $n$. Fix
\[
1\le b\le k,\qquad t=\lceil k/b\rceil,\qquad m=bt\le\rank(A).
\]
The relative clustered gap is
\begin{equation}\label{eq:gap}
\Delta=\min_{1\le i\le m-b}\frac{\lambda_i-\lambda_{i+b}}{\lambda_i}>0.
\end{equation}
When $m=b$, use the empty-minimum convention $\Delta=1$. All logarithms are natural. All assertions concern exact arithmetic and the original, unperturbed matrix $A$.

For a standard real Gaussian $G\in\R^{n\times b}$, define
\begin{equation}\label{eq:algorithm}
\K_q(M,G)=\range[G,MG,\ldots,M^{q-1}G],\qquad
\ahat=Z[Z^TA]_k,
\end{equation}
where $Z$ is any orthonormal basis of this space and $[\cdot]_k$ is a truncated SVD. These are precisely the algorithmic conventions of RA-04~\cite{repo}.

\begin{theorem}[RA-04]\label{thm:main}
There is an absolute constant $C>0$ such that, for every input above and every
$0<\varepsilon,\delta<1/2$, choosing
\begin{equation}\label{eq:target}
q=\left\lceil\frac{C}{\sqrt\varepsilon}
\left(t\log\frac2\Delta+\log\frac{n}{\delta\varepsilon}\right)\right\rceil
\end{equation}
in \eqref{eq:algorithm} gives, with probability at least $1-\delta$,
\begin{align}
\nn{A-\ahat}_\xi&\le(1+\varepsilon)\nn{A-[A]_k}_\xi,
&&\xi\in\{2,F\},\label{eq:errors}\\
\big|\nn{Av_i}_2^2-\lambda_i\big|&\le\varepsilon\lambda_{k+1},
&&1\le i\le k,\label{eq:right}
\end{align}
where $v_i$ are the ordered top $k$ \emph{right} singular vectors of $\ahat$.
The zero-optimal-error case is included.
\end{theorem}

The earlier partial reports proved an all-input estimate with $t\log(2m/\Delta)$ in place of $t\log(2/\Delta)$. They left a dimension-polynomial, gap-exponential interpolation bound unproved. The present argument proves that bound directly. It does not rely on the earlier restricted-width cluster theorem or on the earlier fractional-moment recurrence.

The two substantive external inputs are explicitly isolated: the real-algebraic conic tube theorem in Section~\ref{sec:probability}, and the standard deterministic Krylov good-start theorem in Section~\ref{sec:algorithm}. The root-covering, contour, descent, and interpolation arguments are proved below. None uses an unproved small-ball or interpolation premise.

\section{Head interpolation and its singular locus}\label{sec:interpolation}
For the moment, work only with positive nodes
$\lambda_1\ge\cdots\ge\lambda_m>0$, where $m=bt$ and \eqref{eq:gap} holds. Write $\Lambda=\diag(\lambda_1,\ldots,\lambda_m)$ and, for any $H\in\R^{m\times b}$, put
\begin{equation}\label{eq:Khead}
K_\lambda(H)=[H,\Lambda H,\ldots,\Lambda^{t-1}H]\in\R^{m\times m}.
\end{equation}
Its rows are $h_i^T[I_b,\lambda_i I_b,\ldots,\lambda_i^{t-1}I_b]$.
Whenever it is invertible, set
\begin{equation}\label{eq:eval}
E_{\lambda,H}(x)=[I_b,xI_b,\ldots,x^{t-1}I_b]K_\lambda(H)^{-1}.
\end{equation}
For $y\in\R^m$, $P(x)=E_{\lambda,H}(x)y$ is the unique vector polynomial of degree at most $t-1$ satisfying $h_i^TP(\lambda_i)=y_i$.

\begin{lemma}[Generic invertibility]\label{lem:generic}
The polynomial $D_\lambda(H)=\det K_\lambda(H)$ is nonzero and homogeneous of degree $m$ in the $mb$ entries of $H$. Consequently a standard Gaussian $H$ gives an invertible $K_\lambda(H)$ almost surely.
\end{lemma}
\begin{proof}
Every entry of $K_\lambda(H)$ is linear in $H$, so its determinant is homogeneous of degree $m$ unless identically zero. To exclude that possibility, set row $i$ of $H$ equal to $e_{1+((i-1)\bmod b)}^T$. Row and column permutations then give $b$ scalar Vandermonde blocks. The nodes in block $r$ are
$\lambda_r,\lambda_{r+b},\ldots,\lambda_{r+(t-1)b}$, which are distinct by \eqref{eq:gap}. Thus all blocks are invertible. For $t=1$, this witness is simply $H=I_b$.
A nonzero real polynomial vanishes on a Lebesgue-null set, proving the Gaussian assertion.
\end{proof}

Define the closed algebraic cone
\begin{equation}\label{eq:sigma}
\Sigma_\lambda=\{H\in\R^{m\times b}:D_\lambda(H)=0\}.
\end{equation}
The distance to this cone is always measured in the Frobenius norm on the \emph{starting matrix} $H$, not in an unrestricted perturbation norm on the raw matrix $K_\lambda(H)$.

\begin{lemma}[Translation invariance]\label{lem:translation}
For every real $x$,
\begin{equation}\label{eq:translation}
\Sigma_{\lambda-x}=\Sigma_\lambda,
\qquad E_{\lambda-x,H}(0)=E_{\lambda,H}(x)
\quad(H\notin\Sigma_\lambda).
\end{equation}
In fact, $D_{\lambda-x}(H)=D_\lambda(H)$ for every $H$.
\end{lemma}
\begin{proof}
Expanding $(\Lambda-xI)^jH$ gives
$K_{\lambda-x}(H)=K_\lambda(H)R_x$, where $R_x$ is block upper triangular with $(\ell,j)$ block
$\binom{j}{\ell}(-x)^{j-\ell}I_b$ for $0\le\ell\le j\le t-1$.
All diagonal blocks are $I_b$, so $\det R_x=1$. This proves the determinant and locus assertions. Translating a vector interpolant as $P(y)\mapsto P(y+x)$ proves the evaluation identity by uniqueness.
\end{proof}

This exact invariance will be important: the probability event used below is unchanged by every common shift of the nodes. A union bound over evaluation points is unnecessary.

\section{A logarithmic root covering}\label{sec:cover}
The next elementary covering has total length proportional to $h$, not to the degree. That distinction prevents a $t\log t$ loss.

\begin{lemma}[One-dimensional root covering]\label{lem:cover}
Let $a_1,\ldots,a_s\in\R$, with multiplicities allowed, and let $0<h\le1$. There is a finite union $V$ of closed intervals, of total length at most $6h$, such that
\begin{equation}\label{eq:rootproduct}
\prod_{\ell=1}^s\min\{1,|y-a_\ell|\}\ge(h/e)^s
\end{equation}
whenever $y$ is outside the interior of $V$.
For $s=0$, take $V=\varnothing$ and the empty product equal to one.
\end{lemma}
\begin{proof}
For $s>0$, let $d_j(y)$ be the $j$th smallest value of $|y-a_\ell|$ and define the open bad set
\[
\Omega=\{y:d_j(y)<jh/s\text{ for at least one }j\in\{1,\ldots,s\}\}.
\]
For each such pair $(y,j)$, the open interval $I(y,j)$ centered at $y$ with radius $jh/s$ contains at least $j$ of the points, counted with multiplicity.

From this family, select an interval of largest available length, remove all intervals intersecting it, and repeat. There are only $s$ possible lengths. Selected intervals are disjoint, and each contains at least one point in its interior, so at most $s$ selections occur. If $j_\ell$ corresponds to a selected interval, disjointness gives $\sum_\ell j_\ell\le s$. Hence the sum of the selected lengths is at most $2h$.

Every discarded interval meets a selected interval at least as long. Its center is therefore at distance less than twice the radius of that selected interval. Consequently $\Omega$ lies inside the interior of the union $V$ of the closed triple dilates of the selected intervals. This union has total length at most $6h$.

Outside $\Omega$, $d_j(y)\ge jh/s$. Since $jh/s\le1$,
\[
\prod_{\ell=1}^s\min\{1,|y-a_\ell|\}
\ge\prod_{j=1}^s\frac{jh}{s}
=h^s\frac{s!}{s^s}\ge(h/e)^s.
\]
The last inequality follows from
$\sum_{j=1}^s\log j\ge\int_1^s\log u\,du\ge s\log s-s$.
Every point outside the interior of $V$ is outside $\Omega$.
\end{proof}

\begin{corollary}[Radial lower bound]\label{cor:radial}
Let $\phi$ be a real or complex scalar polynomial of degree $s$ with $\phi(0)=1$. For $0<h\le1$ there is a union $V\subset\R$ of closed intervals of total length at most $6h$ such that
\begin{equation}\label{eq:radial}
|\phi(z)|\ge\left(\frac{h}{2e}\right)^s
\quad\text{whenever }\log|z|\notin\operatorname{int}V.
\end{equation}
\end{corollary}
\begin{proof}
The assertion is immediate for $s=0$. Otherwise write
$\phi(z)=\prod_{\ell=1}^s(1-z/\zeta_\ell)$, where every $\zeta_\ell\ne0$ because $\phi(0)=1$. Apply Lemma~\ref{lem:cover} to $a_\ell=\log|\zeta_\ell|$. If $y=\log|z|$, the reverse triangle inequality gives
\[
|1-z/\zeta_\ell|\ge|1-e^{y-a_\ell}|
\ge\tfrac12\min\{1,|y-a_\ell|\}.
\]
For a nonnegative exponent, use $e^u-1\ge u$; for a negative exponent, use
$1-e^{-u}\ge(1-e^{-1})\min\{1,u\}\ge\tfrac12\min\{1,u\}$.
Multiplication and \eqref{eq:rootproduct} prove the result.
\end{proof}

\section{The annular moment estimate}\label{sec:annular}
For a row set $J$, $H_J$ denotes the corresponding submatrix. If $|J|\le b$, the notation $\sm(H_J)$ means its smallest row singular value, equivalently
\[
\nn{H_J^Tz}_2\ge\sm(H_J)\nn z_2\quad(z\in\R^{|J|}).
\]
If every contiguous $b$-row submatrix has smallest singular value at least $\mu$, then the same holds for every smaller contiguous row set: extend that set to $b$ contiguous rows and pad $z$ with zeros.

\begin{lemma}[Paired nondegeneracy]\label{lem:paired}
Suppose $H\in\R^{m\times b}$ satisfies $\nn H_2\le R$ and every contiguous $b$-row submatrix has smallest singular value at least $\mu>0$. Let
$u\in\R^m$ satisfy
\begin{equation}\label{eq:moments}
\nn u_2=1,\qquad H^T\Lambda^j u=0\quad(1\le j\le t-1).
\end{equation}
For every scalar polynomial $\phi$ of degree at most $t-1$ with $\phi(0)=1$,
\begin{equation}\label{eq:paired}
\nn{\diag(\phi(\lambda_1),\ldots,\phi(\lambda_m))u}_2
\ge\frac{\mu\gamma\rho}{4\sqrt m\,R},
\qquad
\gamma=\frac{\Delta}{64m},\quad
\rho=\left(\frac{\Delta}{128e}\right)^{t-1}.
\end{equation}
No invertibility of $K_\lambda(H)$ is assumed in this lemma.
\end{lemma}
\begin{proof}
Use Corollary~\ref{cor:radial} with $h=\Delta/64$. Since $\deg\phi\le t-1$, the resulting lower bound on its allowed circles is at least $\rho$. Enlarge its covering set to
\begin{equation}\label{eq:U}
U=V\ \cup\ \bigcup_{i=1}^m[\log\lambda_i-\gamma,\log\lambda_i+\gamma].
\end{equation}
The total length of $U$ is at most
\[
6\frac{\Delta}{64}+2m\frac{\Delta}{64m}=\frac{\Delta}{8}.
\]
Every connected component therefore has length at most $\Delta/8$. It contains at most $b$ nodes $\log\lambda_i$, counting multiplicities. Indeed, $b+1$ such nodes would include some $i$ and $i+b$ and give
\[
\log\lambda_i-\log\lambda_{i+b}\le\Delta/8,
\]
whereas \eqref{eq:gap} gives this difference at least $-\log(1-\Delta)\ge\Delta$. When $t=1$, there are only $b$ nodes in total, so this conclusion is immediate.

Consider a component $[a,c]$ containing at least one node. Its node indices form a contiguous set $J$ of size at most $b$. Both endpoints are outside the interior of $V$, and their distance from every $\log\lambda_i$ is at least $\gamma$: an endpoint at a smaller distance would be interior to a node interval in \eqref{eq:U}. Thus, on each of the circles $|z|=e^a$ and $|z|=e^c$,
\begin{equation}\label{eq:contourbounds}
|\phi(z)|\ge\rho,
\qquad
\max_i\frac{\lambda_i}{|z-\lambda_i|}
\le\frac1{1-e^{-\gamma}}\le\frac2\gamma.
\end{equation}
The last inequality uses $0<\gamma\le1$.

Put $w=\diag(\phi(\lambda_i))u$. For any $z$ outside the nodes, the function of $x$
\[
r_z(x)=x\frac{\phi(x)-\phi(z)}{x-z}
\]
is a polynomial of degree at most $t-1$ with zero constant coefficient. The moment equations \eqref{eq:moments} imply
\begin{equation}\label{eq:resolvent}
\phi(z)H^T\Lambda(zI-\Lambda)^{-1}u
=H^T\Lambda(zI-\Lambda)^{-1}w.
\end{equation}
Let $\Gamma_J$ consist of the outer circle $|z|=e^c$ traversed counterclockwise and the inner circle $|z|=e^a$ traversed clockwise. Divide \eqref{eq:resolvent} by $z\phi(z)$ and integrate:
\begin{equation}\label{eq:annularidentity}
H_J^Tu_J
=\frac1{2\pi i}\int_{\Gamma_J}
\frac{H^T\Lambda(zI-\Lambda)^{-1}w}{z\phi(z)}\,dz.
\end{equation}
To verify the left side, integrate the equivalent left integrand
$H^T\Lambda(zI-\Lambda)^{-1}u/z$. For each coordinate,
\[
\frac{\lambda_i}{z(z-\lambda_i)}
=\frac1{z-\lambda_i}-\frac1z.
\]
The two circles enclose exactly the head nodes indexed by $J$ in their annulus; their contributions from zero cancel. Repeated nodes give the sum of their row contributions. Zeros of $\phi$ inside the annulus cause no difficulty: the identity is used on the boundary, and residues are evaluated only for the left integrand, which contains no $1/\phi$.

On either circle, its length cancels the factor $1/|z|$. Consequently \eqref{eq:contourbounds} and \eqref{eq:annularidentity} give
\[
\nn{H_J^Tu_J}_2\le\frac{4R}{\gamma\rho}\nn w_2,
\qquad
\nn{u_J}_2\le\frac{4R}{\mu\gamma\rho}\nn w_2.
\]
The components containing nodes partition all $m$ indices, and there are at most $m$ of them. Summing the squared inequalities yields
\[
1=\nn u_2^2\le m\left(\frac{4R}{\mu\gamma\rho}\right)^2\nn w_2^2,
\]
which proves \eqref{eq:paired}.
\end{proof}

The proof uses at most $b$ rows at a time, but does not fix a partition into $b$-sized spectral clusters. It therefore allows arbitrary nonzero-width clusters, exact repeated values, and unaligned spectra subject only to \eqref{eq:gap}.

\section{Extrapolation, descent, and distance to singularity}\label{sec:descent}
For $H\notin\Sigma_\lambda$, define
\begin{equation}\label{eq:f}
f_\lambda(H)=\frac1{\nn{E_{\lambda,H}(0)}_2}.
\end{equation}
This is positive and continuous on the open set $\R^{m\times b}\setminus\Sigma_\lambda$.

\begin{lemma}[Variational formula and a descent direction]\label{lem:descent}
For $H\notin\Sigma_\lambda$,
\begin{equation}\label{eq:variational}
f_\lambda(H)=
\min_{\substack{P\in\R^b[x],\ \deg P\le t-1\\\nn{P(0)}_2=1}}
\left(\sum_{i=1}^m[h_i^TP(\lambda_i)]^2\right)^{1/2}.
\end{equation}
The minimum is attained. If $H$ has the $\mu,R$ bounds of Lemma~\ref{lem:paired}, there is a matrix $B\in\R^{m\times b}$ with $\nn B_F=1$ such that
\begin{equation}\label{eq:descent}
\limsup_{s\downarrow0}\frac{f_\lambda(H+sB)-f_\lambda(H)}s\le-\alpha,
\qquad\alpha=\frac{\mu\gamma\rho}{4\sqrt m\,R}.
\end{equation}
\end{lemma}
\begin{proof}
The invertible linear map $K_\lambda(H)$ identifies coefficient vectors with data $y_i=h_i^TP(\lambda_i)$; the constant coefficient is $E_{\lambda,H}(0)y$. Thus the minimum in \eqref{eq:variational} is
$\min\{\nn y_2:\nn{E_{\lambda,H}(0)y}_2=1\}=1/\nn{E_{\lambda,H}(0)}_2$. A right singular vector for the largest singular value attains this minimum. In particular, there is no compactness assumption on an unbounded polynomial coefficient set.

Let $P$ attain the minimum and set
\[
y_i=h_i^TP(\lambda_i),\qquad f=f_\lambda(H),\qquad u=y/f.
\]
Varying the positive-degree coefficient vectors of $P$ freely leaves $P(0)$ fixed. First-order optimality of the squared residual gives
$H^T\Lambda^ju=0$ for $1\le j\le t-1$.
Put $a=P(0)$ and $\phi(x)=a^TP(x)$. Then $\nn a_2=1$, $\phi(0)=1$, and $|\phi(\lambda_i)|\le\nn{P(\lambda_i)}_2$.
Lemma~\ref{lem:paired} therefore implies
\begin{equation}\label{eq:W}
D:=\left(\sum_i u_i^2\nn{P(\lambda_i)}_2^2\right)^{1/2}\ge\alpha.
\end{equation}
Let row $i$ of $B$ be $-u_iP(\lambda_i)^T/D$. Then $\nn B_F=1$. Keep this particular $P$ fixed when replacing $H$ by $H+sB$. It remains feasible for \eqref{eq:variational}; its residual norm is differentiable at $s=0$ because $f>0$. Its derivative is
\[
\sum_i u_i\,B_i^TP(\lambda_i)=-D\le-\alpha.
\]
Since $f_\lambda(H+sB)$ is no greater than this feasible residual norm and equality holds at $s=0$, \eqref{eq:descent} follows.
This argument does not require a unique minimizing polynomial or a simple largest singular value of $E(0)$.
\end{proof}

\begin{lemma}[A local error bound]\label{lem:distance}
Fix $H_0\notin\Sigma_\lambda$ and $r_0>0$. Suppose every matrix $X$ in the closed Frobenius ball $\overline B_F(H_0,r_0)$ has spectral norm at most $R$ and contiguous $b$-row singular values at least $\mu$. With $\alpha$ as in \eqref{eq:descent},
\begin{equation}\label{eq:distance}
f_\lambda(H_0)\ge\frac\alpha2
\min\{r_0,\dist_F(H_0,\Sigma_\lambda)\}.
\end{equation}
\end{lemma}
\begin{proof}
Suppose the strict reverse inequality holds, write $f_0=f_\lambda(H_0)$, and put $r=2f_0/\alpha$. Then $r<r_0$ and $r<\dist_F(H_0,\Sigma_\lambda)$. The function
\[
g(X)=f_\lambda(X)+\frac\alpha2\nn{X-H_0}_F
\]
is continuous on the compact closed radius-$r$ ball, which is disjoint from $\Sigma_\lambda$. On its boundary, $g(X)=f_\lambda(X)+f_0>f_0=g(H_0)$. Its minimum is therefore attained in the interior, at some $X_*$. Lemma~\ref{lem:descent} gives a unit direction with upper directional derivative of $f_\lambda$ at most $-\alpha$. The distance penalty has directional increment at most $\alpha/2$ times the step, by the triangle inequality. Thus $g$ has a strictly decreasing direction at $X_*$, a contradiction.
\end{proof}

Lemma~\ref{lem:distance} is the bridge between an extrapolation norm and algebraic geometry. It does not assert that a raw Krylov singular value equals distance to the starting-matrix singular locus. Such an identification would generally be false.

\section{A probability bound uniform in the head spectrum}\label{sec:probability}
We use the following consequence of a published real-algebraic tube theorem. The ambient space is real, as is the Gaussian starting block.

\begin{lemma}[Imported conic tube bound]\label{lem:tube}
Let $\Sigma\ne\{0\}$ be the zero set of a nonzero homogeneous real polynomial of degree $d_0\ge1$ on $\R^{p+1}$, $p\ge1$. For a standard Gaussian $g\in\R^{p+1}$, define
$\mathscr C_\Sigma(g)=\nn g_2/\dist(g,\Sigma)$. Then, for $v\ge3d_0p$,
\begin{equation}\label{eq:tube}
\Prob\{\mathscr C_\Sigma(g)\ge v\}\le\frac{40d_0p}{v}.
\end{equation}
\end{lemma}
\begin{proof}[Derivation from the imported theorem]
Theorem~1.1 of B\"urgisser, Cucker, and Lotz~\cite{bcl}, in its average-case specialization, bounds the probability by
\begin{equation}\label{eq:bclfull}
4\sum_{j=1}^{p-1}\binom pj(2d_0)^j(1+v^{-1})^{p-j}v^{-j}
+\frac{2pO_p}{O_{p-1}}(2d_0/v)^p,
\end{equation}
where $O_p$ is the surface area of $S^p$. Their statement permits singular real algebraic sets. The Gaussian direction is uniform on $S^p$, and the condition number is homogeneous of degree zero, so the spherical estimate applies unchanged.

The sum in \eqref{eq:bclfull} is at most
$4[\exp(3d_0p/v)-1]\le24d_0p/v$, since $v\ge3d_0p$ and $e^z-1\le2z$ for $0\le z\le1$.
Also $O_p/O_{p-1}=\int_0^\pi\sin^{p-1}\theta\,d\theta<4$ and $2d_0/v\le1$, so the last term is at most $16d_0p/v$. Adding proves \eqref{eq:tube}.
Only \eqref{eq:bclfull} is imported; this simplification is included to make the constants explicit.
\end{proof}

For completeness, the elementary Gaussian estimates needed with it are proved here.

\begin{lemma}[Gaussian row-block estimates]\label{lem:gauss}
If $g\sim N(0,1)$, $H\in\R^{m\times b}$ is standard Gaussian, and $S\in\R^{b\times b}$ is standard Gaussian, then for $a>0$,
\begin{equation}\label{eq:gauss}
\Prob\{|g|<a\}\le a,\qquad \E\nn H_F^2=mb,
\qquad \Prob\{\sm(S)<a\}\le b^{3/2}a.
\end{equation}
\end{lemma}
\begin{proof}
The scalar bound follows by integrating the normal density, which is at most $1/\sqrt{2\pi}$, over an interval of length $2a$. The expectation is the sum of the entry variances.
For the last bound, let $d_i$ be the distance from row $i$ of $S$ to the span of its other rows. Almost surely $S$ is invertible and the norm of column $i$ of $S^{-1}$ is $d_i^{-1}$. Hence
\[
\nn{S^{-1}}_2\le\nn{S^{-1}}_F\le\sqrt b\max_i d_i^{-1}.
\]
Conditioned on the other rows, $d_i$ has distribution $|g|$. If $\sm(S)<a$, at least one $d_i<\sqrt b\,a$. The union bound proves the claim. No independence among the row-distance events is required.
\end{proof}

\begin{theorem}[Anchored extrapolation]\label{thm:zero}
For every admissible positive head spectrum, standard Gaussian $H\in\R^{m\times b}$, and $0<\eta<1/2$,
\begin{equation}\label{eq:zero}
\Prob\left\{\nn{E_{\lambda,H}(0)}_2
\le 2^{24}m^8\eta^{-4}\left(\frac{128e}{\Delta}\right)^t\right\}\ge1-\eta.
\end{equation}
For $t\ge2$, the bound holds on an explicitly specified event depending on the spectrum only through $\Sigma_\lambda$ and the stated lower gap bound $\Delta$.
\end{theorem}
\begin{proof}
First take $t\ge2$, so $m\ge2$ and $N=mb\ge2$. Put $p=N-1$ and
\[
R_0=2\sqrt{N/\eta},\qquad
\tau=\frac{\eta}{4mb^{3/2}}.
\]
Consider the simultaneous event
\begin{align}
\nn H_F&\le R_0,\label{eq:event1}\\
\sm(H_{\{r,\ldots,r+b-1\}})&\ge\tau\quad(1\le r\le m-b+1),\label{eq:event2}\\
\nn H_F&\ge\eta/4,\label{eq:event3}\\
\frac{\nn H_F}{\dist_F(H,\Sigma_\lambda)}&\le\frac{160mp}{\eta}.\label{eq:event4}
\end{align}
Each failure probability is at most $\eta/4$. For \eqref{eq:event1}, use Markov's inequality and $\E\nn H_F^2=N$. For \eqref{eq:event2}, apply \eqref{eq:gauss} to at most $m$ contiguous $b$-row blocks. For \eqref{eq:event3}, the event of a smaller Frobenius norm implies $|H_{11}|<\eta/4$. For \eqref{eq:event4}, use Lemma~\ref{lem:tube} with $d_0=m$ and Lemma~\ref{lem:generic}. The cone contains nonzero matrices with a zero row, and it is not the whole ambient space. Thus every hypothesis of the imported theorem holds. A union bound gives probability at least $1-\eta$.

On this event,
\begin{equation}\label{eq:distgood}
\dist_F(H,\Sigma_\lambda)\ge\frac{\eta^2}{640mp}
\ge\frac{\eta^2}{640m^3}=:d_*.
\end{equation}
Here $p\le mb\le m^2$. In the closed ball of radius $r_0=\tau/2$ around $H$, perturbation of a smallest singular value by at most the perturbation norm gives row-block lower bound
\[
\mu=\tau/2\ge\frac{\eta}{8m^{5/2}}.
\]
Every matrix in the ball has spectral norm at most $R=2R_0\le4m/\sqrt\eta$; indeed $r_0<R_0$. Therefore Lemmas~\ref{lem:paired} and \ref{lem:distance} apply throughout this ball, with
\begin{equation}\label{eq:alphagood}
\alpha=\frac{\mu\gamma\rho}{4\sqrt m\,R}
\ge\frac{\eta^{3/2}\Delta\rho}{8192m^5},
\qquad \rho=\left(\frac{\Delta}{128e}\right)^{t-1}.
\end{equation}
Also $r_0\ge\eta/(8m^{5/2})\ge d_*$. Consequently
\begin{align}
f_\lambda(H)&\ge\frac\alpha2d_*
\ge\frac{\eta^{7/2}\Delta\rho}{10485760m^8}\nonumber\\
&\ge\frac{\eta^4}{2^{24}m^8}\left(\frac{\Delta}{128e}\right)^t.
\label{eq:fbound}
\end{align}
The last step only weakens the bound: $10485760<2^{24}$, $\eta^{7/2}\ge\eta^4$, and $\Delta\rho\ge(\Delta/(128e))^t$. Inverting gives \eqref{eq:zero}.

When $t=1$, $m=b$ and $E_{\lambda,H}(0)=H^{-1}$, independently of the nodes. Lemma~\ref{lem:gauss} gives $\nn{H^{-1}}_2\le m^{3/2}/\eta$ with probability at least $1-\eta$. This is smaller than the right-hand side of \eqref{eq:zero}, with $\Delta=1$. It also handles the one-dimensional case $m=b=1$, for which the conic tube theorem was not invoked.
\end{proof}

\begin{theorem}[Uniform tail interpolation]\label{thm:uniform}
Under the same assumptions,
\begin{equation}\label{eq:uniform}
\Prob\left\{\sup_{0\le x\le\lambda_m}\nn{E_{\lambda,H}(x)}_2
\le 2^{24}m^8\eta^{-4}\left(\frac{128e}{\Delta}\right)^t\right\}\ge1-\eta.
\end{equation}
\end{theorem}
\begin{proof}
For $t\ge2$, use the single event \eqref{eq:event1}--\eqref{eq:event4} for the original spectrum. For every $0\le x<\lambda_m$, all shifted nodes $\lambda_i-x$ are positive and
\[
\frac{(\lambda_i-x)-(\lambda_{i+b}-x)}{\lambda_i-x}
\ge\frac{\lambda_i-\lambda_{i+b}}{\lambda_i}\ge\Delta.
\]
By Lemma~\ref{lem:translation}, their singular locus is \emph{exactly} $\Sigma_\lambda$. Thus the same four events, local ball, and estimates \eqref{eq:distgood}--\eqref{eq:fbound} apply to every shifted spectrum with the same constants. This is a deterministic simultaneous implication on that event, not an uncountable union of probability claims.

The evaluation identity in \eqref{eq:translation} proves the bound for every $x<\lambda_m$. Since $E_{\lambda,H}(x)$ is a matrix polynomial, continuity gives it also at $x=\lambda_m$. For $t=1$, $E(x)=H^{-1}$ is constant and the assertion follows from Theorem~\ref{thm:zero}.
\end{proof}

This establishes the sufficient interpolation estimate left open in the earlier reports. The dimension and failure probability appear only as fixed polynomial factors; their logarithms are not multiplied by $t$.

\section{Transfer to the original randomized algorithm}\label{sec:algorithm}
We now return to $M=AA^T$ and the original $G$. Take a fixed orthogonal eigenbasis $U=[U_m\ U_\perp]$ and write
\[
U^TMU=\diag(\Lambda,\Lambda_\perp),\qquad
U^TG=\begin{bmatrix}H\\T\end{bmatrix}.
\]
Orthogonal invariance gives standard Gaussian entries in both blocks. No eigenvectors are selected using $G$.

Let
\[
B=[G,MG,\ldots,M^{t-1}G],\qquad
K_T=[T,\Lambda_\perp T,\ldots,\Lambda_\perp^{t-1}T].
\]
Almost surely,
\begin{equation}\label{eq:graph}
BK_\lambda(H)^{-1}=U\begin{bmatrix}I_m\\F\end{bmatrix},
\qquad F=K_TK_\lambda(H)^{-1}.
\end{equation}
At a tail node $x_j\in[0,\lambda_m]$, row $j$ of $F$ is $T_j^TE_{\lambda,H}(x_j)$. Hence
\begin{equation}\label{eq:graphnorm}
\nn F_2^2\le\nn F_F^2
\le\nn T_F^2\sup_{0\le x\le\lambda_m}\nn{E_{\lambda,H}(x)}_2^2.
\end{equation}

Apply Theorem~\ref{thm:uniform} with $\eta=\delta/2$, and intersect its event with
$\nn T_F^2\le2nb/\delta$. The latter fails with probability at most $\delta/2$ by Markov's inequality, since $\E\nn T_F^2=(n-m)b\le nb$. No independence is needed for this intersection. We obtain, with probability at least $1-\delta$,
\begin{equation}\label{eq:L}
L:=1+\nn F_2^2
\le1+2^{57}nbm^{16}\delta^{-9}\left(\frac{128e}{\Delta}\right)^{2t}
\le1+2^{57}n^{18}\delta^{-9}\left(\frac{128e}{\Delta}\right)^{2t}.
\end{equation}

For clarity, the graph-to-overlap reduction is proved rather than assumed. Let $J_k\in\R^{m\times k}$ contain the first $k$ coordinate vectors and let $F_k=FJ_k$. Then
\[
Q=U\begin{bmatrix}J_k\\F_k\end{bmatrix}(I_k+F_k^TF_k)^{-1/2}
\]
has orthonormal columns in $\range(B)$, and
\begin{equation}\label{eq:good}
\nn{(U_k^TQ)^{-1}}_2^2
=\nn{I_k+F_k^TF_k}_2\le L.
\end{equation}
This handles $k$ not divisible by $b$ directly.

\begin{theorem}[Imported deterministic good-start convergence]\label{thm:imported}
There is an absolute $C_0\ge1$ with the following property. If the column space of a starting matrix $B$ contains orthonormal $Q\in\R^{n\times k}$ satisfying \eqref{eq:good}, then projection onto $\K_s(M,B)$ followed by a rank-$k$ SVD truncation satisfies \eqref{eq:errors}--\eqref{eq:right} for sufficiently many steps,
\[
s\ge\max\left\{2,\left\lceil\frac{C_0}{\sqrt\varepsilon}\log\frac{nL}{\varepsilon}\right\rceil\right\}.
\]
\end{theorem}

This is Imported Theorem~3.2 in Chen, Epperly, Meyer, Musco, and Rao~\cite{chen}, with its sufficient-iteration $O(\cdot)$ notation expressed using an absolute constant. Their Problem~1.1 explicitly specifies \emph{right} singular vectors. The theorem invokes the good-start analysis in Appendix~F of Meyer, Musco, and Musco~\cite{mmm}. It is the only convergence theorem imported here. The case $\lambda_{k+1}=0$ is also proved directly below rather than relying on a zero-denominator limit in that analysis.

\begin{proof}[Proof of Theorem~\ref{thm:main}]
Suppose first $\lambda_{k+1}>0$. On the probability-$1-\delta$ event above, \eqref{eq:L} and \eqref{eq:good} hold. Direct multiplication of the block powers gives the exact identity
\begin{equation}\label{eq:nested}
\K_s(M,B)=\K_{s+t-1}(M,G).
\end{equation}
Put $X=t\log(2/\Delta)$ and $Y=\log(n/(\delta\varepsilon))$. The bound \eqref{eq:L} implies
\begin{align*}
\log\frac{nL}{\varepsilon}
&\le58\log2+19\log n+9\log(1/\delta)+\log(1/\varepsilon)
 +2t\log(128e/\Delta)\\
&\le 80(X+Y).
\end{align*}
For the last inequality, use $58\log2\le58X$,
$\log(128e/\Delta)\le9\log(2/\Delta)$, and
$19\log n+9\log(1/\delta)+\log(1/\varepsilon)\le19Y$.
Consequently a choice such as $C=80C_0+4$ in \eqref{eq:target} ensures
$q-t+1\ge\max\{2,\lceil C_0\log(nL/\varepsilon)/\sqrt\varepsilon\rceil\}$.
Theorem~\ref{thm:imported} and \eqref{eq:nested} prove all three guarantees for the actual output of \eqref{eq:algorithm}.

If $\lambda_{k+1}=0$, then $\rank(A)=k$ and the hypothesis $m\le\rank(A)$ forces $m=k$, so $b$ divides $k$. In the chosen eigenbasis,
\[
U^T[MG,M^2G,\ldots,M^tG]
=\begin{bmatrix}\Lambda K_\lambda(H)\\0\end{bmatrix}.
\]
Both head factors are invertible almost surely. Thus $\K_{t+1}(M,G)$ contains $\range(A)$, giving $ZZ^TA=A$ and an exact best rank-$k$ approximation. The right singular-vector energies are exactly $\lambda_i$. The above choice of $C$ ensures $q\ge t+1$, so this case is included as well.
\end{proof}

\section{Scope, failure modes avoided, and review obligations}\label{sec:audit}
The proof is uniform in all of $n,d,k,b,t,\varepsilon,\delta$ and in the positive head spectrum. It places no bound on $\lambda_1/\lambda_m$ and requires no positive consecutive gap. It leaves the matrix $A$ and Gaussian distribution in RA-04 unchanged. The word ``perturbation'' in the cited tube theorem does not denote an extra algorithmic perturbation here: only its average-case spherical specialization is used.

Several potential shortcuts are deliberately not used. The determinant being nonzero does not alone give a quantitative bound; the descent lemma and the tube theorem supply that missing step. Smallest singular values of raw monomial Krylov matrices are not asserted to be gap-only. Noncommuting matrix-valued Lagrange products are not introduced. Finally, the uniform evaluation bound follows from a common event and exact translation invariance, not from an uncountable union bound.

\begin{longtable}{p{0.30\textwidth}p{0.64\textwidth}}
\toprule
Proof obligation & Disposition\\
\midrule
\endhead
Repeated head values & They occur at most $b$ times under the gap assumption; the residue-class witness and annular row groups allow repetitions.\\
Unaligned or wide clusters & Only the $b$-step relative gap and contiguous row blocks are used. No fixed cluster centers or width condition occur.\\
Roots at head nodes & Contour boundaries avoid all root moduli in the covering. Roots inside a contour do not invalidate the boundary identity.\\
Complex roots & The reverse triangle inequality gives a bound on an entire circle from the root modulus.\\
Multiple extremal singular values & The descent proof freezes one minimizer and uses an upper directional derivative, not differentiability of the minimum.\\
Singular algebraic locus & The imported tube theorem covers singular real algebraic sets. No smoothness of the determinant-zero set is imposed.\\
Probability dependence & Four elementary/conic events are combined by a union bound. Their independence is not asserted.\\
All tail evaluations & The determinant-zero locus is invariant under common node translations. One event controls all $x<\lambda_m$; continuity gives the endpoint.\\
Nondivisible target rank & The first $k$ graph columns give the overlap certificate directly.\\
Zero tail error & Exact range recovery at depth $t+1$ is proved separately.\\
Right-vector energy guarantee & Imported Theorem~3.2 is used with the right-vector convention of its Problem~1.1, not a substituted left-vector estimate.\\
Universal constants & A fixed polynomial prefactor and exponent-$t$ gap factor give \eqref{eq:L}; $C=80C_0+4$ is one sufficient choice relative to the imported constant.\\
\bottomrule
\end{longtable}

\textbf{Evidence level.} The complete argument passed a separate Codex AI-agent informal mathematical audit. Under the repository's policy this supports \emph{Solved}~\cite{guidelines}. The independent report accompanies this submission. No external human peer review, formal proof certificate, Lean verification or novelty-priority claim is asserted.

The independent audit examined the annular identity \eqref{eq:annularidentity}, the passage from a frozen minimizing polynomial to the descent inequality \eqref{eq:descent}, and the compact-ball argument in Lemma~\ref{lem:distance}. These are the new steps connecting spectral clustering to the published algebraic probability theorem. Their assumptions and proofs are stated in full above; they are not outstanding conjectural premises.

\section{Reproducible checks}\label{sec:checks}
The companion scripts are tests of formulas and implementations, not substitutes for the universal arguments. Rational fixtures are expressly not presented as Gaussian samples. Floating-point observations are not probability certificates.

\textbf{Exact checks.} The rational-arithmetic suite passes five residue-class rank witnesses, including two $t=1$ endpoints, and three interpolation fixtures. It verifies six vector resolvent identities, 105 positive-pole coordinate residues, translation and determinant invariance, Schur-complement formulas for the anchored minimum, and direct graph identities. Four exact interval-cover fixtures include repeated logarithmic root locations and verify the total-length bound. Additional checks verify 13 symbolic Pascal translation identities and a tied-singular-value fixture in which the frozen-minimizer descent remains valid without uniqueness.

\textbf{Head-matrix diagnostics.} Across 96 fixtures with $b,t\in\{1,2,3,4\}$, geometric, repeated, and nonzero-width head spectra were tested. The largest relative discrepancy between the analytic singular-value derivative and the frozen-minimizer derivative was $6.05\cdot10^{-14}$; the corresponding finite-difference discrepancy was $6.75\cdot10^{-5}$. The maximum normalized moment residual was $4.03\cdot10^{-16}$, and the maximum relative translation residual was $8.10\cdot10^{-13}$. The stated interpolation bound exceeded every sampled norm, but its conservatism and these finite observations do not establish a probability bound.

\textbf{Algorithmic diagnostics.} Two-pass block Arnoldi with numerical-rank detection in binary64 produced 81 records across seven spectra and three seeds per spectrum. Every record directly measures both norm errors and the ordered right-singular-vector energy defect. At $\varepsilon=0.1$, 43 records satisfied all tested criteria. The table gives the smallest \emph{tested} depth at which all three seeds passed. Failed smaller-depth records are retained. Neither the table nor its success counts estimate or certify the universal constant in Theorem~\ref{thm:main}.
\begin{center}
\begin{tabular}{lrrrr}
\toprule
Fixture & $n$ & $b$ & $k$ & Tested depth\\
\midrule
Exactly repeated clusters & 72 & 4 & 24 & 18\\
Wide, unaligned spectrum & 64 & 4 & 19 & 10\\
Mixed multiplicities at the boundary & 42 & 3 & 10 & 9\\
Nonzero-width clusters & 45 & 3 & 14 & 10\\
Scalar starting block & 20 & 1 & 5 & 12\\
Full-size starting block & 32 & 6 & 6 & 4\\
Zero optimal error & 40 & 3 & 12 & 5\\
\bottomrule
\end{tabular}
\end{center}
For the zero-error fixture, ``passed'' means that the spectral residual, Frobenius residual, and energy defect are each at most $10^{-10}$; the theorem is an exact-arithmetic assertion, not this tolerance test. Its exact recovery depth is $t+1=5$.

\textbf{Spectral-scale stress tests.} Six further interpolation solves used 180-decimal-digit arithmetic. For the scalar fixture with head spectrum $(1,\zeta,\zeta/2)$ and $\zeta=10^{-4},10^{-16},10^{-50}$, the clustered gap stayed at $1/2$. The raw smallest singular values were approximately $1.65\cdot10^{-5}$, $1.65\cdot10^{-17}$, and $1.65\cdot10^{-51}$, while the anchored evaluation norms were approximately $4.07350$, $4.07329$, and $4.07329$. This finite diagnostic illustrates the difference between raw conditioning and extrapolation; it does not prove a limiting statement. Normalized inverse residuals are recorded as backward residual diagnostics, not forward-error certificates.

The recorded environment is Python 3.13.5, NumPy 2.3.5, SymPy 1.14.0, and mpmath 1.3.0. Platform-dependent floating-point variation is expected.


From the package root, run:
\begin{Verbatim}[fontsize=\small]
python -m pip install -r requirements.txt
python tests/exact_checks.py --output results/exact_checks.json
python tests/numerical_checks.py --output-dir results
./build.sh
\end{Verbatim}
The tests retain seeds, dimensions, numerical residuals, and failed smaller-depth approximation criteria. Identity-check failures stop execution. The source and a machine-readable source inventory are included. The prior packages are identified by hashes in the provenance file; no earlier unproved interpolation estimate is assumed by this proof.

\clearpage
\begin{thebibliography}{9}
\bibitem{repo}
A. Townsend, \emph{Open Problems in Numerical Linear Algebra}, entry RA-04, ``Clustered singular-value gaps in randomized block Krylov approximation.'' Repository statement consulted September 12, 2026 (US Eastern time).
\url{https://github.com/ajt60gaibb/OpenProblemsInNLA/tree/main/randomized-and-low-rank-approximation/RA-04}.

\bibitem{chen}
T. Chen, E. N. Epperly, R. A. Meyer, C. Musco, and A. Rao,
\emph{Does block size matter in randomized block Krylov low-rank approximation?}
SODA 2026, pp. 1026--1046.
\href{https://doi.org/10.1137/1.9781611978971.42}{doi:10.1137/1.9781611978971.42}.
Version used: \href{https://arxiv.org/abs/2508.06486v2}{arXiv:2508.06486v2}, Problem~1.1 and Imported Theorem~3.2, PDF pp. 1 and 6. The PDF title page identifies Christopher Musco and Akash Rao.

\bibitem{mmm}
R. A. Meyer, C. Musco, and C. Musco,
\emph{On the Unreasonable Effectiveness of Single Vector Krylov Methods for Low-Rank Approximation}, SODA 2024, pp. 811--845.
\href{https://arxiv.org/abs/2305.02535}{arXiv:2305.02535}.
Appendix~F is the convergence-transfer provenance cited by~\cite{chen}.

\bibitem{bcl}
P. B\"urgisser, F. Cucker, and M. Lotz,
\emph{The probability that a slightly perturbed numerical analysis problem is difficult}, Mathematics of Computation 77 (2008), 1559--1583.
\href{https://doi.org/10.1090/S0025-5718-08-02060-7}{doi:10.1090/S0025-5718-08-02060-7}.
Theorem~1.1 in the inspected preprint \href{https://arxiv.org/abs/math/0610270v1}{arXiv:math/0610270v1}, titled \emph{The probability that a small perturbation of a numerical analysis problem is difficult}, PDF p. 4. The real theorem, not the separate complex conic theorem, is used.

\bibitem{guidelines}
\emph{Open Problems in Numerical Linear Algebra}, README, ``Problem status,'' and CONTRIBUTING.md. Consulted September 12, 2026 (US Eastern time).
\url{https://github.com/ajt60gaibb/OpenProblemsInNLA}.
\end{thebibliography}
\end{document}
